% 欠定线性回归中初始化的零空间分量被梯度下降原样保留。
% 一个训练观测 x=(1,1)^T、y=1，半平方损失 (w1+w2-1)^2/2，步长 eta=1/4。
% 残差 r=w1+w2-1 每步乘 1-2eta=1/2；更新 w <- w - eta r (1,1)^T。
% 零初始化：0 -> (.25,.25) -> (.375,.375) -> (.4375,.4375) -> B=(.5,.5)。
% 初值 w0=(-0.6,0)：行空间分量 (-0.3,-0.3)，零空间分量 v=(-0.3,0.3)；
%   迭代 (-0.2,0.4) -> (0,0.6) -> (0.1,0.7) -> 极限 B+v=(0.2,0.8)。
% 坐标：x=1mm,y=1mm，参数每单位 40 mm。
\begingroup
\begin{tikzpicture}[figure picture,x=1mm,y=1mm]
  \def\s{40}
  % 坐标轴
  \draw[-{Stealth[length=1.6mm]},FigureStroke,line width=.6pt] (-0.85*\s,0)--(1.3*\s,0) node[right]{$w_1$};
  \draw[-{Stealth[length=1.6mm]},FigureStroke,line width=.6pt] (0,-0.25*\s)--(0,1.3*\s) node[above]{$w_2$};
  \node[figure note,anchor=north east] at (0,0) {$0$};
  % 行空间方向 w1=w2 与零空间方向 w1=-w2
  \draw[FigureMuted,line width=.4pt,dash pattern=on 1.5pt off 1.5pt] (-0.2*\s,-0.2*\s)--(1.15*\s,1.15*\s);
  \node[figure note,anchor=west] at (1.0*\s,1.12*\s) {行空间 $w_1=w_2$};
  \draw[FigureMuted,line width=.4pt,dash pattern=on 1.5pt off 1.5pt] (0.2*\s,-0.2*\s)--(-0.62*\s,0.62*\s);
  \node[figure note,anchor=east] at (-0.62*\s,0.66*\s) {零空间 $w_1=-w_2$};
  % 插值直线
  \draw[FigureBlue,line width=1.1pt] (-0.25*\s,1.25*\s)--(1.25*\s,-0.25*\s);
  \node[font=\figurefont\footnotesize,text=FigureBlue,anchor=west] at (1.04*\s,-0.06*\s) {插值解 $w_1+w_2=1$};
  % 零初始化轨迹
  \draw[FigureRed,line width=.9pt,-{Stealth[length=1.5mm]}] (0,0)--(0.25*\s,0.25*\s);
  \draw[FigureRed,line width=.9pt,-{Stealth[length=1.5mm]}] (0.25*\s,0.25*\s)--(0.375*\s,0.375*\s);
  \draw[FigureRed,line width=.9pt] (0.375*\s,0.375*\s)--(0.5*\s,0.5*\s);
  \foreach \p in {0.25,0.375,0.4375}{\fill[FigureRed] (\p*\s,\p*\s) circle (0.7);}
  \fill[FigureRed] (0.5*\s,0.5*\s) circle (1.3);
  \node[font=\figurefont\footnotesize,text=FigureRed,anchor=north west] at (0.56*\s,0.40*\s) {$\widehat{\bm w}=(0.5,0.5)$};
  % 非零初值轨迹
  \fill[FigureYellow] (-0.6*\s,0) circle (1.1);
  \node[font=\figurefont\footnotesize,anchor=north] at (-0.6*\s,-1.5) {$\bm w^{(0)}=(-0.6,0)$};
  \draw[FigureYellow!80!black,line width=.9pt,-{Stealth[length=1.5mm]}] (-0.6*\s,0)--(-0.2*\s,0.4*\s);
  \draw[FigureYellow!80!black,line width=.9pt,-{Stealth[length=1.5mm]}] (-0.2*\s,0.4*\s)--(0,0.6*\s);
  \draw[FigureYellow!80!black,line width=.9pt] (0,0.6*\s)--(0.2*\s,0.8*\s);
  \foreach \a/\b in {-0.2/0.4,0/0.6,0.1/0.7}{\fill[FigureYellow!80!black] (\a*\s,\b*\s) circle (0.7);}
  \fill[FigureYellow!80!black] (0.2*\s,0.8*\s) circle (1.3);
  \node[font=\figurefont\footnotesize,anchor=south west] at (0.22*\s,0.82*\s) {$\widehat{\bm w}+\bm v=(0.2,0.8)$};
  % 零空间分量 v
  \draw[FigureInk,line width=.6pt,-{Stealth[length=1.5mm]}] (0.5*\s,0.5*\s)--(0.2*\s,0.8*\s);
  \node[figure edge label] at (0.38*\s,0.62*\s) {$\bm v$};
  \draw[FigureInk,line width=.6pt,-{Stealth[length=1.5mm]}] (0,0)--(-0.3*\s,0.3*\s);
  \node[figure edge label,anchor=north east] at (-0.15*\s,0.13*\s) {$\bm v$};
\end{tikzpicture}
\endgroup
